\section{Initial Subsystem Requirements}
\subsection{Mechanical}
These subsystem requirements can be applied to all concepts:
\begin{itemize}
    \item \textbf{MechR1:} The system shall be compatible with the launch system and be able to be carried by the rocket.
    \item \textbf{MechR2:} The system shall withstand \textit{g} forces at launch of up to 6\textit{g}.
    \item \textbf{MechR3:} The system shall survive extreme vibrations encountered at launch.
    \item \textbf{MechR4:} The system shall retain all subsystem from \textit{g} forces up to 6\textit{g} at launch condition.
    \item \textbf{MechR5:} The system shall be made up of materials that can withstand variation in environmental temperature of $-230^{\circ}$ to $120^{\circ}$. 
    \begin{itemize}
        \item This extreme range is provided to take into account of the different concepts operating in LEO and deep space (L2).
    \end{itemize}
    \item \textbf{MechR6:} The chassis shall not deform past an acceptable degree due to temperature fluctuation.
    \item \textbf{MechR7:} The material of the chassis shall not impede communication signals between the satellite and ground station.
    \item \textbf{MechR8:} The chassis shall be able to withstand and protect the internal subsystem from collisions caused by space debris of size up to $1$ mm large and travelling at around 10 km/s \cite{ESA2020HubbleDebris}.
    \item \textbf{MechR9:} The chassis shall be able to withstand thermal cracking caused by the shrinking and expansion of material over subsequent temperature cycles.
    \item \textbf{MechR10:} The chassis shall not obstruct any optical and observation capabilities of the payload and other subsystem.
    \item \textbf{MechR11:} Solar panels shall not have a failure rate of more than 0.137\% \cite{stewart2021final}.
    \item \textbf{MechR12:} Solar panels shall have a range of rotation of $10^{\circ}$. 
\end{itemize}
\textbf{MechR1,2,3,4} are required in order for the satellite to be compatible with current generation launch systems as well as survive the launch process.

\textbf{MechR5,6,8,9} are required for the general structural strenght of the satellite, aiding its durability.

\textbf{MechR7,10,11,12} are to ensure the telescope can conduct its mission successfully.
\subsection{GNC}
\begin{itemize}
    \item \textbf{GNCR1:} The system shall be able to track and identify its position in the orbit at any phase of the orbit.
    \begin{itemize}
        \item The satellite shall have sensors that can track its position in any phase of the orbit (e.g If the satellite is in a solar eclipse, it must have sensors other than a solar sensor)
    \end{itemize}
    \item \textbf{GNCR2:} The satellite shall be able to change its orientation along its axes. 
    \item \textbf{GNCR3:} The satellite shall be able to deduce its true anomaly at any point in the orbit.
    \item \textbf{GNCR4:} The system shall keep the telescope pointed at a specific point with no more than 0.007 arcseconds of deviation.
    \item \textbf{GNCR5:} The system shall keep the telescope's accuracy to within 60 arcseconds.
    \item \textbf{GNCR6:} The system shall be able to rotate the telescope at an angular speed of at least $0.00175$ rad/s.
\end{itemize}
\textbf{GNCR4, GNCR5, GNCR6} were derived from HST. Our concept must at the very least meet these requirements \cite{HubbleADC}.

\textbf{GNCR1, GNCR2, GNCR3} are general requirements for the basis of an attitude determination and control system (ADC).
\subsection{Propulsion}
\begin{itemize}
    \item \textbf{PR1:} The system shall use propellants and propulsion systems available within the scheduled plan for launch, within the next 5 years. 
    \item \textbf{PR2:} Propellants ejected from the system shall not interfere with any optics and observation sensors in the satellite. 
    \item \textbf{PR3:} The propulsion system shall produce enough thrust to perform any required manoeuvres by the satellite.
    \item \textbf{PR4:} The propulsion system shall have sufficient propellant available to maintain orbit for long periods. This will ensure that the satellite is functional for at least 15 years. 
    \item \textbf{PR5:} The propulsion system shall not cause excess vibrations that may damage the satellite.
    \item \textbf{PR6:} The propulsion system must not interfere with the imaging system, ensuring that no contamination results from any thrust manoeuvres. 
\end{itemize}
\subsection{Payload}

\begin{itemize}
    \item \textbf{PyldR1:} The sensor system shall detect both metallic and non-metallic debris under 10 cm in diameter to meet mission objectives. This will be verified via ground testing.
    \item \textbf{PyldR2:} Instruments shall withstand launch vibrations (20 Hz–10,000 Hz) and sound waves (up to 220dB) without damage. This will be validated by analysis and testing.
    \item \textbf{PyldR3:} Sensors shall detect debris without physical contact, preserving instrument integrity.
    \item \textbf{PyldR4:} The sensor system shall determine debris position, velocity (approx. 7.8 km/s), and last known orbit, integral for mission tracking.
    \item \textbf{PyldR5:} The system shall have a consistent power supply ranging from 200-800 watts is essential for continuous sensor operation, confirmed by simulation.
    \item \textbf{PyldR6:} The camera sensor shall be able to cover a wide spectrum from infrared (IR) to visible light (V) and ultraviolet (UV).
\end{itemize}

\begin{itemize}
    \item \textbf{PyldR7:} The system shall show structural resilience against extreme temperatures (–230°C to +125°C) and radiation levels (100–1000 rad/year).
    \item \textbf{PyldR8:} The payload structure shall withstand launch-induced vibrations and sound waves to operate reliably in orbit.
    \item \textbf{PyldR9:} The propulsion system shall accurately maintain the satellite's trajectory for effective data collection, validated through design simulations.
\end{itemize}

\textbf{PyldR1, PyldR3, PyldR4} are essential for accurate debris detection and tracking, ensuring mission objectives are met.

\textbf{PyldR2, PyldR7, PyldR8} are necessary to withstand harsh launch and operational conditions, ensuring structural integrity.

\textbf{PyldR5, PyldR6} are critical for continuous operation and broad spectral imaging capabilities.

\subsection{Electrical Subsystem Requirements}

\begin{itemize}

    \item \textbf{ElecR1:}
    The subsystem shall include a component that generates at least 110\% of the satellite’s weekly power requirement to support continuous operation of all subsystems, including Electrical, Propulsion, Mechanical, and Payload, over each Earth week.
    
    \item \textbf{ElecR2:}
    The subsystem shall include a component with the capacity to store at least 110\% of the total electrical energy requirement to account for operational variability. 

    \item \textbf{ElecR3:}
    The subsystem shall include a Power Management Unit (PMU) to regulate and distribute electrical power to each component at appropriate voltage and current levels. 

    \item \textbf{ElecR4:}
    The PMU shall incorporate protection mechanisms, including over-voltage, over-current, and thermal protection, to prevent damage to onboard systems. 
   
    \item \textbf{ElecR5:}
    Redundant power pathways shall be included to allow the subsystem to continue functioning in the event of a fault in the primary power path. The PMU and PHM shall monitor and manage these pathways to maintain system resilience.
    \begin{itemize}
        \item ElecR1 to R5 reflects the requirement of UR2 - The need of satliate to be self-sustainable. Shall be achieved with power generation, stroage capcity and PMU unit.
    \end{itemize}
    \item \textbf{ElecR6:}
    The subsystem shall be capable of transmitting and receiving telecommunication signals with ground control on Earth at intervals no longer than 10 minutes.

    \item \textbf{ElecR7:}
    The subsystem shall have an onboard flash memory unit capable of storing command, image, and telemetry data for up to two weeks during transmission loss.

    \item \textbf{ElecR8:}
    The subsystem shall include a Data Processing Unit (DPU) that encodes and decodes data transmissions, including command, telemetry, and payload data encryption and decryption.
     \begin{itemize}
        \item ElecR6 to R8 reflects the requirement of UR2 - The need of bidirectional communication. 
    \end{itemize}
    \item \textbf{ElecR9:}
    The DPU shall support routing and receiving electronic signals across all signal-using components, functioning as a data bus for efficient data exchange.

    \item \textbf{ElecR10:}
    The subsystem shall include a location measurement unit comprising a 3 Degrees of Freedom (DOF) Inertial Measurement Unit (IMU), gyroscope, accelerometer, and magnetometer to provide current location information (for space applications).

    \item \textbf{ElecR11:}
    The subsystem’s onboard flight computer shall receive data remotely and the local navigation system and send electronic control signals to the propulsion and mechanical systems (for space applications).
    
    \item \textbf{ElecR12:}
    The flight computer shall compute predictive trajectory adjustments and generate control signals for mechanical components to adjust trajectory at a frequency of at least 1000 Hz (for space applications).
    
    \item \textbf{ElecR13:}
    The flight computer shall be capable of making autonomous decisions for trajectory adjustments in emergency collision avoidance scenarios (for space applications).
     \begin{itemize}
        \item ElecR9 to R13 reflects the requirement of UR3 - The need of satellite to be navigational. 
    \end{itemize}

    \item \textbf{ElecR14:}
    The subsystem shall ensure that any interrupted transmission due to environmental challenges can be re-established promptly after conditions improve (for space applications).

    \item \textbf{ElecR15:}
    The system must include a temperature regulation unit with passive or active thermal control to maintain operational temperatures across all subsystems.

    \item \textbf{ElecR16:}
    All components in the subsystem shall be radiation-hardened to withstand cumulative radiation exposure consistent with the satellite’s orbital environment (for space applications).

    \item \textbf{ElecR17:}
    The subsystem shall support communication and other critical subsystems to withstand environmental disturbances, including extreme vibrations, temperature variations, and radiation levels.
     \begin{itemize}
        \item ElecR14 to R17 reflects the requirement of UR7 - The need of receiving stable images. It furthers the redundancy of the whole system by predicting failure. 
    \end{itemize}

    
\end{itemize}



\subsection{Operations and Concepts of Operation}
\begin{itemize}
    \item \textbf{OpR1:} The onboard system of any satellite shall frequently transmit captured data back to a ground station. This approach minimises the need for large onboard data storage, thereby reducing the satellite's overall mass and, consequently, the fuel required for operation.
    \item \textbf{OpR2:} The system shall be capable of reorienting, maintaining dynamic stability, and pointing at any target region in space, provided it's not obstructed by a celestial body. Orientation control is crucial for maintaining a direct line of communication with the ground station for data transmission. Stability ensures that onboard cameras can capture a region of space over prolonged periods without perturbation, enabling detailed imaging.
    \item \textbf{OpR3:} The satellite shall be capable of maneuvering to avoid collisions with debris. This will ensure that the lifetime of the satellite is not compromised.
    \item \textbf{OpR4:} The satellite shall remain in the desired orbit for its full lifetime to ensure ideal operations.
    \item \textbf{OpR5:} The satellite shall remain on a stable orbit to reduce the impact of orbit decay.
    \item \textbf{OpR6:} The satellite shall be able to operate semi-autonomously when there is a temporary communication drop, without the continuous need for human intervention.
    \begin{itemize}
        \item There is a lag time between communications hence for unexpected events, the satellite should be able to respond to emergencies
        \item Data continuity during blackouts is important. The satellite should be able to store and continue with other tasks without first having communicated previous data.
        \item However, enough onboard storage shall be provided to hold a minimum of 24 hours’ worth of observation data in case of communication delays or errors.
    \end{itemize}
\end{itemize}